\chapter{ベクトルのノルム}
    \section{諸定理}
        \subsection{\texorpdfstring{$\norm{\bm{x}}_{p+a} \leq \norm{\bm{x}}_p\;(\bm{x}\in\realNumbers^n,\;p\geq 1,\;a\geq 0)$}{||x||\_\{p+a\} ≦ ||x||\_p (x∈R\string^n, p≧1, a≧0)}}
            \begin{proof}
                \quad\par
                $a=0$のときは明らか。$a>0$のときを考える。
                $x_1,\dotsc,x_n$のうち非零要素が1つ以下のときは明らか。
                非零要素が2つ以上のときを考える。
                一般性を失わずに$|x_1|,\dotsc,|x_r|>0\;(2\leq r\leq n)$とする。
                $\derivLong{\norm{\bm{x}}_p}{p}{} < 0$ を示せば良い。
                \begin{align*}
                    \derivLong{\norm{\bm{x}}_p}{p}{} < 0 &\iff \derivLong{\log \norm{\bm{x}}_p}{p}{} < 0 \iff -\frac{1}{p^2}\log \sum_{i=1}^r |x_i|^p + \frac{1}{p}\times\frac{\sum_{i=1}^r |x_i|^p\log |x_i|}{\sum_{i=1}^r |x_i|^p} < 0 \\
                    &\iff \frac{\sum_{i=1}^r |x_i|^p\log |x_i|}{\sum_{i=1}^r |x_i|^p} < \frac{1}{p}\log \sum_{i=1}^r |x_i|^p \\
                    &\iff \sum_{i=1}^r w_i\log |x_i|^p < \log M \quad \left( M \coloneqq \sum_{i=1}^r |x_i|^p,\; w_i \coloneqq |x_i|^p/M \right) \\
                    &\iff \sum_{i=1}^r w_i\log M w_i < \log M \iff \sum_{i=1}^r w_i(\log M + \log w_i) < \log M \\
                    &\iff \sum_{i=1}^r w_i\log w_i < 0 \quad \left( \because\;\sum_{i=1}^r w_i = 1 \right) \quad = \text{true} \quad (\because\; 0<w_i<1)
                \end{align*}
            \end{proof}
        \subsection{\texorpdfstring{$1\leq p\leq q,\;\bm{x}\in\complexNumbers^n$}{1≦p≦q, x∈C\string^n}のとき\texorpdfstring{$\norm{\bm{x}}_p \leq n^{1/p-1/q}\norm{\bm{x}}_q$}{||x||\_p ≦ n\string^(1/p-1/q)||x||\_q}}
            \begin{proof}
                \quad\par
                $\bm{a},\bm{b}\in\complexNumbers^n$とする。$r=q/p, s=q/(q-p)$とすると$1\leq r,s$および$1/r+1/s=1$を満たすのでH\"olderの不等式より次式が成り立つ。
                \[ \sum_{i=1}^n|a_ib_i| \leq \left(\sum_{i=1}^n|a_i|^r\right)^\frac{1}{r}\left(\sum_{i=1}^n|b_i|^s\right)^\frac{1}{s} \]
                $|a_i| = |x_i|^p,\;b_i=1$とすると次式を得る。
                \[ \sum_{i=1}^n|x_i|^p \leq \left(\sum_{i=1}^n|x_i|^q\right)^\frac{p}{q}n^{1-p/q} \]
                両辺の$p$乗根をとって次式を得る。
                \[ \norm{\bm{x}}_p \leq n^{1/p-1/q}\norm{\bm{x}}_q \]
            \end{proof}
        \subsection{和のノルムとノルムの和が等しい時}
            \begin{shadebox}
                $\vx,\vy\in\field^n$が非零であるとき$\norm{\vx+\vy}=\norm{\vx}+\norm{\vy} \iff ^\exists k>0 \st \vy=k\vx$
            \end{shadebox}
            \begin{proof}
                \quad\par
                $\Leftarrow$は明らか。$\Rightarrow$を示す。
                \[\norm{\vx+\vy}=\norm{\vx}+\norm{\vy}\]
                の両辺を2乗して整理すると
                \begin{equation}
                    \inprod{\vx}{\vy} = \norm{\vx}\norm{\vy}\
                \end{equation}
                よって
                \[\inprod{\vx}{\vy}^2-\norm{\vx}^2\norm{\vy}^2=0\]
                ところでこれは$s$に関する二次方程式$\norm{s\vx+\vy}^2=0$の判別式の値が$0$であることを意味する。
                よって重解が存在し、それを$-k$とおくと
                \begin{equation}
                    \vy = k\vx
                \end{equation}
                式(1)に代入して
                \[k\norm{\vx}^2 = |k|\norm{\vx}^2\]
                $\vx\neq\bm{0}$だから$\norm{\vx}>0$なので$k=|k|\quad\therefore k\geq0$。
                さらに$\vx,\vy\neq\bm{0}$であったから式(2)より$k\neq0$。結局$k>0$
            \end{proof}
        \subsection{和のノルムとノルムの和が等しい時(2)}
            \begin{shadebox}
                $\bm{x}_1,\dotsc,\bm{x}_m\in\field^n$に対して$\norm{\sum_{i=1}^{m}{\bm{x}_i}}=\sum_{i=1}^{m}{\norm{\bm{x}_i}}$であるとき、任意の$k,l\in\{1,2,\dotsc,m\}$に対して$\norm{\bm{x}_k+\bm{x}_l}=\norm{\bm{x}_k}+\norm{\bm{x}_l}$
            \end{shadebox}
            \begin{proof}
                \quad\par
                $k=l$については明らか。$k\neq l$について示す。
                三角不等式より
                \[\norm{\sum_{i=1}^{m}{\bm{x}_i}} \leq \norm{\bm{x}_k+\bm{x}_l} + \sum_{i\neq k,l}{\norm{\bm{x}_i}} \leq \sum_{i=1}^{m}{\norm{\bm{x}_i}}\]
                仮定より最右辺が$\norm{\sum_{i=1}^{m}{\bm{x}_i}}$と等しいから結局
                \[\norm{\sum_{i=1}^{m}{\bm{x}_i}} \leq \norm{\bm{x}_k+\bm{x}_l} + \sum_{i\neq k,l}{\norm{\bm{x}_i}} \leq \norm{\sum_{i=1}^{m}{\bm{x}_i}}\]
                \[\therefore\; \norm{\sum_{i=1}^{m}{\bm{x}_i}} = \norm{\bm{x}_k+\bm{x}_l} + \sum_{i\neq k,l}{\norm{\bm{x}_i}} = \sum_{i=1}^{m}{\norm{\bm{x}_i}}\]
                \[\therefore\; \norm{\bm{x}_k+\bm{x}_l} = \norm{\bm{x}_k}+\norm{\bm{x}_l}\]
            \end{proof}
        % \subsection{$\bm{x}\in\realNumbers^n,\norm{\bm{x}}_2=1\;\Rightarrow\;\sum_{i=1}^n{|x_i|}\leq\sqrt{n}$}
        % 	\label{th:単位球面上の点の座標の絶対値の総和の最大値}
        % 	\begin{proof}
        % 		\quad\par
        % 		全要素が正の場合について証明すれば十分である。
        % 		Lagrangeの未定乗数法を使う。
        % 		目的関数を$f(\bm{x})\coloneqq\sum_{i=1}^n{x_i}$、制約条件を$g(\bm{x})\coloneqq\norm{\bm{x}}_2^2-1=0$とし、Lagrange関数を$L(\bm{x},\lambda)\coloneqq f(\bm{x})+\lambda g(\bm{x})$とする。
        % 		最適解を$\bm{x}^*$とすると$\nabla_{\bm{x}}L(\bm{x},\lambda)=\bm{0}$より$\bm{1}_n+2\lambda\bm{x}^*=\bm{0}\;\therefore\;\bm{x}=-\frac{1}{2\lambda}\bm{1}_n$。
        % 		よって$\bm{x}^*$は全要素が等しい。これと$\norm{\bm{x}^*}=1,\bm{x}^*\geq\bm{0}$より$\bm{x}^* = \frac{1}{\sqrt{n}}\bm{1}_n$であり、$f(\bm{x}^*)=\sqrt{n}$
        % 	\end{proof}
        \subsection{\texorpdfstring{$\bm{x}\in\complexNumbers^n,\norm{\bm{x}}_2=1\;\Rightarrow\;\sum_{i=1}^n{|x_i|}\leq\sqrt{n}$}{x∈C\string^n, ||x||\_2=1, ⇒ |x\_1|+...+|x\_n| ≦ √n}}
            \label{th:単位球面上の点の座標の絶対値の総和の最大値}
            \begin{proof}
                \quad\par
                $x_k = r_k e^{i\theta_k},\; (r_k \geq 0, \theta_k \in \realNumbers,\;k=1,\dotsc,n)$とすると$\norm{\bm{x}}_2=1 \iff \sum_{k=1}^n {r_k}^2 = 1,\;\sum_{k=1}^n{|x_k|} = \sum_{k=1}^n r_k$であるから、制約条件$g(\bm{r}) \coloneqq \sum_{k=1}^n {r_k}^2 - 1 = 0,\;\bm{r} \geq \bm{0}_n$の下で関数$f(\bm{r}) \coloneqq \sum_{k=1}^n r_k$の最大値が$\sqrt{n}$以下であること(*)を示せばよい。
                \par
                数学的帰納法で示す。$n=1$のときは明らかに(*)が成り立つ。以下$n=m\geq2$とし、$n=1,\dotsc,m-1$について(*)が成り立つと仮定する。
                $r_1,\dotsc,r_n$の少なくとも1つが0のとき、すなわち制約集合の境界上では変数の個数が$n-1$以下の場合に帰着するから帰納法の仮定より成り立つ。
                \par
                そこで$r_1,\dotsc,r_n$が全て正であるとする。
                Lagrangeの未定乗数法を使う。
                Lagrange関数を$L(\bm{r},\lambda)\coloneqq f(\bm{r})+\lambda g(\bm{r})$とする。
                $\nabla_\bm{r}g(\bm{r}) = 2\bm{r} > \bm{0}_n$であるから、今関心のある領域に特異点は存在しない。
                $r_n = \sqrt{1 - \sum_{k=1}^{n-1} {r_k}^2}$と表されるから、$f(\bm{r}) = \sum_{k=1}^{n-1} r_k + \sqrt{1 - \sum_{k=1}^{n-1} {r_k}^2}$であり、これは$r_1,\dotsc,r_{n-1}$の関数として上に凸な関数である(Hessianを計算すれば対角成分が全て負の対角行列になる)。
                よってLagrangeの未定乗数法で極値を見つければそれは極大値であり、また最大値でもある。
                極値を与える点を$\overset{\circ}{\bm{r}}$とすると$\nabla_{\bm{r}}L(\overset{\circ}{\bm{r}},\lambda)=\bm{0}$より$\bm{1}_n + 2\lambda\overset{\circ}{\bm{r}}=\bm{0}\;\therefore\;\overset{\circ}{\bm{r}}=-\frac{1}{2\lambda}\bm{1}_n$。
                よって$\overset{\circ}{\bm{r}}$は全要素が等しい。
                これと制約条件より$\overset{\circ}{\bm{r}} = \frac{1}{\sqrt{n}}\bm{1}_n$であり、$f(\overset{\circ}{\bm{r}}) = \sqrt{n}$
            \end{proof}
    \section{応用}
        \subsection{超平面と点の距離公式の直感的説明}
            超平面$\bm{w}^\top\bm{x}+b=0$と点$\bm{y}$の距離公式$|\bm{w}^\top\bm{y}+b|/\norm{\bm{w}}$について直感的な説明を考えてみる。
            $-\bm{w}^\top\bm{p} = b$なる$\bm{p}$を考えると$\bm{w}^\top\bm{x}+b=0 \iff \bm{w}^\top(\bm{x}-\bm{p})=0$であり、$\bm{p}$は超平面上の点で、$\bm{w}$は法線ベクトルである。
            よって$|\bm{w}^\top\bm{y}+b| = |\bm{w}^\top(\bm{y}-\bm{p})|$は$\bm{y}$から超平面に降ろした垂線の長さの$\norm{\bm{w}}$倍を表しており、これを$\norm{\bm{w}}$で割れば公式を得る。